\subsection{Class and Interface Declarations}\label{sec:ClassAndInterfaceDeclarations}
Each Java source file\footnote{… or \bdq{ordinary compilation unit}\index{compilation unit}} contains at least one top level class or interface declaration; in most cases, it is just one. And it is strongly discouraged that it would be more than one!

Only one of these top level classes or interfaces can be \lstinline|public|, the others have to be package-local (meaning the modifiers \lstinline|private| and \lstinline|protected| are not allowed). If you really have more than one class in a single source file, the \lstinline|public| one should be the first entry, all others will follow, sorted alphabetically by their names.

With the JPMS\index{JPMS|see {Java Platform Module System}}\index{Java Platform Module System} (Java Platform Module System) the need for non-public top-level classes got nearly obsolete and they should be avoided.

If you need classes or interfaces that are closer related to a top-level class, you should consider to define these as \textit{inner} classes or interfaces. Inner classes can be \lstinline|public|, \lstinline|private| or \lstinline|protected|. An inner class is part of the top level class, but when defined as \lstinline|static|, it can be used independently from the enclosing class.

At least since Java~5, there are not only classes and interfaces that are declared in a Java source file. Today, we have

\begin{itemize}[nosep]
\item{\hyperref[subsec:Class]{classes} (keyword \lstinline|class|)}
\item{\hyperref[subsec:Interface]{interfaces} (keyword \lstinline|interface|)}
\item{\hyperref[subsec:Enum]{enums} (keyword \lstinline|enum|)}
\item{\hyperref[subsec:Record]{records} (keyword \lstinline|record|)}
\item{\hyperref[subsec:Annotation]{annotations} (keyword \lstinline|@interface|)}
\end{itemize}

Technically, records and enums are special types of classes, while an annotation\index{annotation} is somehow a special kind of an interface.

Each of these allows different sections in their declarations; the following tables describe the various sections and the order they should appear in. Each section has a header comment, and inside each section, the elements are ordered alphabetically.

\subsubsection{Class}\label{subsec:Class}
\keeplisting{34}
A skeleton for a Java \lstinline|class| may look like this:
\begin{lstlisting}[numbers=left,caption={Class Skeleton}]
public class MyClass 
{
        /*---------------*\
    ====** Inner Classes **==========================================
        \*---------------*/
    …    
        
        /*-----------*\
    ====** Constants **==============================================
        \*-----------*/
    …
        
        /*------------*\
    ====** Attributes **=============================================
        \*------------*/
    …
            
        /*------------------------*\
    ====** Static Initialisations **=================================
        \*------------------------*/
    …
        
        /*--------------*\
    ====** Constructors **===========================================
        \*--------------*/
    …
        
        /*---------*\
    ====** Methods **================================================
        \*---------*/
    …            
}
//  class MyClass
\end{lstlisting}
 
\begin{filecontents}{ClassParts.tbl}
  \begin{longtable}{|l|X|}
  \caption{Parts of a \lstinline|class| declaration} \\
  \hline 
  Part & Description \\ 
  \hline
  \endfirsthead
  \multicolumn{2}{c}%c
  {\tablename\ \thetable\ -- \textit{Continued from previous page}} \\
  \hline 
  Part & Description \\ 
  \hline
  \endhead
  \multicolumn{2}{r}{\textit{Continued on next page}} \\ 
  \endfoot
  \endlastfoot
  Inner Classes (optional) & Given the class defines inner classes, interfaces, enums or records, either \lstinline|private| or \lstinline|public|, they will come first. Internally, they will follow the same rules as for a top-level class. \\ 
  \hline 
  Constants (optional) & This part takes all constants defined by this class (meaning \lstinline|public static final| fields). All values for the constants are known at compile time, otherwise you should place the respective field to the \textit{Static Initialisations} part. \\ 
  \hline 
  Attributes (optional) & Here go the attributes for the class; technically, this part is in fact optional for a class, but usually a class without any attributes does not make much sense. Attributes are \lstinline|private| (under some exceptional circumstances, they can be \lstinline|protected|) and they can be \lstinline|static|, although this should be avoided. \\ 
  \hline 
  Static Initialisations (optional) & If a class declares \lstinline|static| fields that are not mere constants or their initialisation is more complex than just assigning a compile time constant, these fields go into this part.
  
  The fields should be initialised in a \lstinline|static {...}| block, not through an immediate initialiser.
  
  This is also the right place for the \lstinline|serialVersionUID|\index{serialVersionUID} of a seria­lisable\index{java.lang!Serializable} class. \\ 
  \hline 
  Constructors (optional) & All the constructors for the class go into this part. Technically, it is optional if the class will have only an empty \lstinline|public| default constructor, but the recommendation is to code this, too. \\ 
  \hline 
  Methods (optional) & All methods of the class are in this part. Again, this part is technically optional, but classes without methods are barely useful.  \\ 
  \hline 
 \end{longtable} 
\end{filecontents}
\LTXtable{\linewidth}{ClassParts.tbl}

\subsubsection{Interface}\label{subsec:Interface}
\keeplisting{20}
A skeleton for an \lstinline|interface| may look like this:

\begin{lstlisting}[numbers=left,caption={Interface Skeleton}]
public interface MyInterface 
{
        /*---------------*\
    ====** Inner Classes **==========================================
        \*---------------*/
    …
        
        /*-----------*\
    ====** Constants **==============================================
        \*-----------*/
    …
        
        /*---------*\
    ====** Methods **================================================
        \*---------*/
    …
                
}
//  interface MyInterface
\end{lstlisting}
 
\begin{filecontents}{InterfaceParts.tbl}
  \begin{longtable}{|l|X|}
  \caption{Parts of an \lstinline|interface| declaration} \\
  \hline 
  Part & Description \\ 
  \hline
  \endfirsthead
  \multicolumn{2}{c}%
  {\tablename\ \thetable\ -- \textit{Continued from previous page}} \\
  \hline 
  Part & Description \\ 
  \hline
  \endhead
  \multicolumn{2}{r}{\textit{Continued on next page}} \\ 
  \endfoot
  \endlastfoot
  Inner Classes (optional) & Defining inner classes for an interface is discouraged, although there are some valid use cases for this. So it could make sense to define an \lstinline|enum| with flag values used by a method declared by the interface. If the interface defines inner classes or interfaces, they have to be \lstinline|public| – classes have to be \lstinline|static|, too, and they will be the first part. Internally, they will follow the same rules as for top-level classes. \\ 
  \hline 
  Constants (optional) & This part takes all constants defined by this interface (meaning \lstinline|public static final| fields). All values for the constants are known at compile time. \\ 
  \hline 
  Methods (optional) & All methods of the interface are in this part. Again, this part is technically optional, it can be omitted for so-called \textit{marker interfaces}.  \\ 
  \hline 
 \end{longtable} 
\end{filecontents}
\LTXtable{\linewidth}{InterfaceParts.tbl}

\subsubsection{enum}\label{subsec:Enum}
\index{enum}An \lstinline|enum| is a special case of a Java class. Usually it only contains the enumeration values, but it can have more components. A skeleton for a \lstinline|enum| class may look like this:

\begin{lstlisting}[numbers=left,caption={enum Skeleton}]
public enum MyEnum 
{
        /*------------------*\
    ====** Enum Definitions **=======================================
        \*------------------*/
    …
        
        /*-----------*\
    ====** Constants **==============================================
        \*-----------*/
    …
        
        /*------------*\
    ====** Attributes **=============================================
        \*------------*/
    …
        
        /*------------------------*\
    ====** Static Initialisations **=================================
        \*------------------------*/
    …
        
        /*--------------*\
    ====** Constructors **===========================================
        \*--------------*/
    …
        
        /*---------*\
    ====** Methods **================================================
        \*---------*/
    …
                
}
//  enum MyEnum
\end{lstlisting}
 
\begin{filecontents}{EnumParts.tbl}
  \begin{longtable}{|l|X|}
  \caption{Parts of an \lstinline|enum| declaration} \\
  \hline 
  Part & Description \\ 
  \hline
  \endfirsthead
  \multicolumn{2}{c}%
  {\tablename\ \thetable\ -- \textit{Continued from previous page}} \\
  \hline 
  Part & Description \\ 
  \hline
  \endhead
  \multicolumn{2}{r}{\textit{Continued on next page}} \\ 
  \endfoot
  \endlastfoot
  Enum Definitions & This part is mandatory. Each \lstinline|enum| is an instance of the enum class. \\ 
  \hline 
  Constants (optional) & This part takes all additional constants defined by this \lstinline|enum| (meaning \lstinline|public static final| fields). All values for the constants are known at compile time, otherwise you should place the respective field to the \textit{Static Initialisations} part. \\ 
  \hline 
  Attributes (optional) & Here goes the attributes for the \lstinline|enum|, if any. Each attribute has to be \lstinline|private final| and will be initialised by the constructor. \\ 
  \hline 
  Static Initialisations (optional) & If an \lstinline|enum| declares \lstinline|static final| fields with an initialisation that is more complex than just assigning a compile time constant, these fields go into this part.
  
  The fields should be initialised in a \lstinline|static {...}| block. \\
  \hline 
  Constructors (optional) & All the constructors for the enum (rarely there is more than one) go into this part, if any. A constructor for an \lstinline|enum| must be \lstinline|private|. \\ 
  \hline 
  Methods (optional) & All methods of the \lstinline|enum| are in this part.  \\ 
  \hline 
 \end{longtable} 
\end{filecontents}
\LTXtable{\linewidth}{EnumParts.tbl}

Technically, an \lstinline|enum| can also define inner classes, but this is strongly discouraged, and I am not aware of any use case for this.

\subsubsection{Record}\label{subsec:Record}
\index{record}Records had been introduced to Java with version~14 and were finalised with version~16 \autocite{ORACLE_DOC_RECORD,ORACLE_DOC_LANGUAGE_SPECIFICATION:RecordClasses,Eckel:JavaRecords}. Like an \lstinline|enum|, a record is a restricted form of a Java class. It’s ideal for \bdq{plain data carrier} classes that contain data not meant to be altered. Usually it has only the most fundamental methods, such as constructors and accessors.

A skeleton for a Java record class may look like this, with all parts being optional:
\begin{lstlisting}[numbers=left,caption={Record Skeleton}]
public record MyRecord( ... ) 
{
        /*-----------*\
    ====** Constants **==============================================
        \*-----------*/
    …
        
        /*------------*\
    ====** Attributes **=============================================
        \*------------*/
    …
        
        /*------------------------*\
    ====** Static Initialisations **=================================
        \*------------------------*/
    …
        
        /*--------------*\
    ====** Constructors **===========================================
        \*--------------*/
    …
        
        /*---------*\
    ====** Methods **================================================
        \*---------*/
    …
                
}
//  record MyRecord
\end{lstlisting}
 
\begin{filecontents}{RecordParts.tbl}
  \begin{longtable}{|l|X|}
  \caption{Parts of a \lstinline|record| declaration} \\
  \hline 
  Part & Description \\ 
  \hline
  \endfirsthead
  \multicolumn{2}{c}%
  {\tablename\ \thetable\ -- \textit{Continued from previous page}} \\
  \hline 
  Part & Description \\ 
  \hline
  \endhead
  \multicolumn{2}{r}{\textit{Continued on next page}} \\ 
  \endfoot
  \endlastfoot
  Constants (optional) & This part takes all constants defined by this record (meaning \lstinline|public static final| fields). All values for the constants are known at compile time, otherwise you should place the respective field to the \textit{Static Initialisations} part. \\ 
  \hline 
  Attributes (optional) & Here go the attributes for the class; this part is optional, because usually, all attributes are defined as arguments to the \lstinline|record| definition. \\ 
  \hline 
  Static Initialisations (optional) & If a record declares \lstinline|static| fields that are not mere constants or their initialisation is more complex than just assigning a compile time constant, these fields go into this part.
  
  The fields should be initialised in a \lstinline|static {...}| block. \\ 
  \hline 
  Constructors (optional) & If the record needs additional constructors, these go into this part. It is optional, because in most cases no additional constructor is need. \\ 
  \hline 
  Methods (optional) & All additional methods of the record are in this part. \\ 
  \hline 
 \end{longtable} 
\end{filecontents}
\LTXtable{\linewidth}{RecordParts.tbl}

Same as \lstinline|enum|s can records have inner classes, but this is similarly discouraged, and there is no use case for it.

\subsubsection{Annotation}\label{subsec:Annotation}
Annotations had been introduced with Java~5 \autocite{ORACLE_DOC_LANGUAGE_SPECIFICATION:AnnotationInterfaces}. A skeleton for an annotation interface may look like this:

\keeplisting{8}
\begin{lstlisting}[numbers=left,caption={Annotation Skeleton}]
public @interface MyAnnotation 
{
        /*------------*\
    ====** Attributes **=============================================
        \*------------*/
    …
}
//  @interface MyAnnotation
\end{lstlisting}
 
\begin{filecontents}{AnnotationParts.tbl}
  \begin{longtable}{|l|X|}
  \caption{Parts of an annotation declaration} \\
  \hline 
  Part & Description \\ 
  \hline
  \endfirsthead
  \multicolumn{2}{c}%
  {\tablename\ \thetable\ -- \textit{Continued from previous page}} \\
  \hline 
  Part & Description \\ 
  \hline
  \endhead
  \multicolumn{2}{r}{\textit{Continued on next page}} \\ 
  \endfoot
  \endlastfoot
  Attributes (optional) & The attributes for the annotation go here, if any. An annotation without attributes would be a marker annotation. \\ 
  \hline 
 \end{longtable} 
\end{filecontents}
\LTXtable{\linewidth}{AnnotationParts.tbl}

\subsubsection{Principles}
\begin{enumerate}[label=P\arabic*.]
\item{Only one public top-level class (interface, record, enum, …) per compilation unit. Avoid non-public top-level classes.}
\item{Group the parts of the class as shown above. Remove the comments for unused sections.}
\end{enumerate}


%==============================================================================
% End of file!!
%==============================================================================
